\chapter*{Abstract}
\addcontentsline{toc}{chapter}{Abstract}

This thesis presents the design, implementation, and evaluation of a software-defined FX correlator for the four-element Nooitgedacht radio telescope array. Radio interferometry relies on precise cross-correlation of voltage signals from multiple antennas to produce visibility measurements, which form the foundation of synthesis imaging. The Nooitgedacht facility lacked a digital correlator, limiting its scientific capabilities to single-dish operations. This work addresses that gap by developing a complete Python-based correlator system implementing the Fourier Transform followed by Cross-multiplication (FX) architecture.

The correlator was implemented using a modular software design that separates data ingestion, channelization, delay compensation, and cross-multiplication into distinct processing stages. The F-engine employs windowed Fast Fourier Transforms to channelize time-domain antenna signals into frequency spectra, while the X-engine computes all baseline cross-correlations through complex multiplication and time integration. Geometric delay compensation accounts for path-length differences between antennas, enabling coherent combination of signals from multiple pointing directions. The system processes dual-polarization signals from four 3.7-metre dishes with configurable spectral resolution up to 4096 channels and supports multiple output formats including NumPy arrays, HDF5, and FITS.

A design science research methodology guided the development process through iterative cycles of design, implementation, and evaluation. The correlator architecture prioritises flexibility and extensibility, leveraging Python's scientific computing ecosystem (NumPy, SciPy) while maintaining pathways for GPU acceleration through libraries like CuPy. Validation was performed using both synthetic test signals with known properties and simulated astronomical observations. Accuracy tests confirmed that channelization produces spectra matching analytical predictions within numerical precision limits, delay compensation correctly accounts for geometric phase shifts, and correlation products preserve expected fringe patterns.

Performance evaluation demonstrated that the system achieves near-real-time processing for modest array sizes, handling four antennas with 256–1024 channels at sample rates up to 2048 Hz on standard computing hardware. For a representative configuration (4 antennas, 256 channels, 1-second integration), the correlator processes 10 seconds of observation data in approximately 0.8 seconds, representing a 12× real-time factor. Memory usage remains tractable at under 2 GB for typical observing modes, and the modular architecture facilitates straightforward scaling to larger channel counts or additional antennas.

The completed system transforms the Nooitgedacht array from four independent telescopes into a functioning interferometer capable of producing calibrated visibility measurements. This enables scientific applications including source imaging, spectral line observations, and polarimetry studies. The software correlator also serves as an educational platform for students to gain practical experience in digital signal processing, radio astronomy, and high-performance computing. By demonstrating that flexible, maintainable correlators can be implemented without specialised FPGA or ASIC hardware, this work contributes to the broader accessibility of radio interferometry technology and complements South Africa's investment in large-scale facilities like MeerKAT and the Square Kilometre Array.